\documentclass[11pt,a4paper]{article}
\usepackage[margin=1in]{geometry}
\usepackage[sort&compress,square,numbers]{natbib}
\usepackage{lipsum,times,graphicx,hyperref,cleveref}
\usepackage[labelfont=bf]{caption}

\title{\vspace{-2cm} {\large Supplementary data}\\[-2mm] 
\textbf{Dispersion of Multiferroic Nanoparticles in a Bent-Core Nematic Liquid Crystal: Experimental and Theoretical Study}}
 
\author{Dhananjoy Mandal, Yiwei Wang, Supreet Kaur, Golam Mohiuddin, Apala Majumdar and Aloka Sinha$^\star$\\ 
\small $^\star$Corresponding author. Email: aloka@physics.iitd.ac.in
}
\date{}

%\renewcommand{\figurename}{Fig.}
\renewcommand{\baselinestretch}{1.33} 
\renewcommand{\thepage}{S\arabic{page}}
\renewcommand{\thetable}{S\arabic{table}}
\renewcommand{\thefigure}{S\arabic{figure}}
\renewcommand{\theequation}{S\arabic{equation}}
%==============================Content=========================
\begin{document}\maketitle

\section*{Small-angle X-ray scattering (SAXS) data}
For conclusive proof disordered cybotactic cluster, X-ray measurements have now been performed for the pure 8-F-OH and doped 8-F-OH + 0.2 $ wt\%$ BiFeO$_3$ systems, maintaining the experimental condition exactly the same for both the systems, to check the disordered cybotactic clustering in doped systems. The small-angle X-ray scattering (SAXS) patterns show that the intensity of the peak increases upon lowering the temperature in both pure and doped 8-F-OH indicating the presence of ordered clusters in the nematic phase (N$_{cyb}$ phase). However, the intensity of the doped 8-F-OH system is comparatively less than that of the pure compound for the same X-ray exposure time i.e., 15 min and at the same temperature. For instance, at 105 $^{\circ}{\rm C}$, the intensity of the doped system is 14 a.u. and the pure compound is 34 a.u. suggesting that the cybotactic clustering is distorted upon dispersion of BiFeO$_3$ nanoparticles ( Figure S1). 

In both pure and doped 8-F-OH compounds, the layer spacing (d-value) and the correlation length ($\xi$) gradually increases in the N$_{cyb}$ phase region (Figure S2) upon lowering the temperature. It is well documented that the increase in the size of the smectic-type clusters results in relatively long correlation lengths. Interestingly, the correlation length values are comparatively lower for the doped system. For instance, at 135 $^{\circ}{\rm C}$, the $\xi$ value for the doped 8-F-OH ($\xi$ = 61.77 $\AA$) is less than the pure 8-F-OH ($\xi$ = 70.49 $\AA$) plausibly due to the deformation of the cybotactic clusters after the dispersion of BiFeO3 nanoparticles. From this SAXS experiments it is quite evident that the BiFeO$_3$ nanoparticle at higher temperature inevitably destroys the cybotactic clusters via interacting with the bent-core LC systems while upon lowering the temperature some regrowth of clusters was observed

\begin{figure}[!htb]
  \centering
    \includegraphics[width=0.49\textwidth]{Pure_SAXS.png}
    \includegraphics[width=0.49\textwidth]{Doped_SAXS.png} 
  \caption{Temperature-dependent XRD patterns in the small-angle region: Intensity vs. 2$\theta$ plots of the (a) pure compound 8-F-OH and (b) doped 8-F-OH + 0.2 $ wt\%$ BiFeO$_3$ system}
\end{figure}

\begin{figure}[!htb]
  \centering
    \includegraphics[width=0.49\textwidth]{Pure_correlation_length.png}
    \includegraphics[width=0.49\textwidth]{Doped_correlation_length.png} 
  \caption{Variation of correlation length ($\xi$) and d-spacing with temperature of (c) pure 8-F-OH and (d) doped 8-F-OH + 0.2 $ wt\%$ BiFeO$_3$ system.}
\end{figure}



\end{document}